\chapter{Représentation des résultats dans le tableau de bord}

\section{Fonction du tableau de bord dans l'architecture}
Le tableau de bord constitue le niveau de synthèse de la couche signal. Alors que
la page Signaux est orientée vers l'analyse d'un titre donné, le tableau de bord
est conçu pour répondre à une autre question : \emph{quels sont, à un horizon
précis, les titres qui méritent une attention prioritaire ?}

La route \texttt{/dashboard} s'appuie sur la vue définie dans
\texttt{frontend/app/v1/page.tsx}. Les données utilisées sont chargées à partir
de fichiers dépendant de l'horizon : \texttt{/data/scores-short.json},
\texttt{/data/scores-medium.json} et \texttt{/data/scores-long.json}. Cette
organisation rappelle une fois encore que la lecture du signal n'est jamais
indépendante de son horizon de construction.

\section{Structure des données consommées}
L'unité de restitution dans le tableau de bord est le titre, modélisé par le type
\texttt{DashboardStock}. Pour chaque action, ce modèle contient un score global,
une étiquette d'interprétation, les contributions par famille et plusieurs
métadonnées descriptives.

Cette transformation est conceptuellement importante. Dans le moteur, l'unité
fondamentale est la variante, puis la famille. Dans le tableau de bord,
l'utilisateur raisonne d'abord par symbole. Le tableau de bord opère donc une
projection de la structure méthodologique vers une structure de décision orientée
métier.

\section{Agrégation globale par titre}
À l'échelle du tableau de bord, le score global d'un titre peut être interprété
comme une synthèse des familles effectivement disponibles. Dans sa forme la plus
simple, cette synthèse peut s'écrire
\begin{equation}
\mathrm{Score}_{global}(a,h)=\frac{1}{N_{a,h}}\sum_{f=1}^{N_{a,h}} \mathrm{Score}^{(\%)}_{f,h}(a),
\label{eq:dashboard-global}
\end{equation}
où \(a\) désigne une action, \(h\) un horizon, et \(N_{a,h}\) le nombre de
familles effectivement prises en compte pour cette action et cet horizon.

L'équation~\ref{eq:dashboard-global} ne doit pas être lue comme une simple
moyenne arithmétique déconnectée de la méthode amont. Chacun des termes
\(\mathrm{Score}^{(\%)}_{f,h}(a)\) est lui-même issu de l'agrégation pondérée
des représentants de la famille, telle qu'elle a été établie au chapitre
précédent. Le score du tableau de bord est donc une synthèse de synthèses,
mais une synthèse qui reste traçable.

\section{Codage visuel et lisibilité}
Le tableau de bord repose sur trois éléments complémentaires.
\begin{itemize}
\item Le tableau principal ordonne les actions selon le score global ou selon une
famille donnée.
\item Les étiquettes textuelles, portées notamment par \texttt{SignalBadge},
traduisent les zones de score en vocabulaire lisible.
\item Les éléments graphiques de type barre de score rendent visible l'intensité
de l'orientation, et pas seulement son signe.
\end{itemize}

La conversion d'un score numérique en libellé suit une grille discrète. Pour le
score global, l'axe \([-100,100]\) est découpé en cinq zones principales :
achat fort, achat, neutre, vente, vente forte. Cette quantification permet une
lecture rapide et homogène d'un grand nombre de lignes, ce qui est précisément la
fonction du tableau de bord.

\section{Lecture multi-horizons}
L'une des forces du tableau de bord réside dans la possibilité de consulter un
même univers de titres selon plusieurs horizons. Ce point mérite d'être souligné
car il prolonge directement le soin méthodologique apporté au calibrage des
familles.

En pratique, un titre peut présenter un signal positif en court terme et une
lecture beaucoup plus prudente en long terme. Une telle divergence n'est pas une
incohérence du système ; elle traduit simplement le fait que les fenêtres
d'évaluation, les grilles paramétriques et la dynamique de marché ne sont pas les
mêmes. Le tableau de bord expose donc une pluralité de lectures temporelles,
plutôt qu'un verdict unique artificiellement unifié.

\section{Continuité entre synthèse et explication}
Chaque ligne du tableau principal constitue également un point d'entrée vers la
page Signaux. Le lien portant le symbole transmet le titre et l'horizon à la
route \texttt{/signals}. Cette continuité est fondamentale :
\[
\text{tableau de bord} \rightarrow \text{page Signaux} \rightarrow \text{détail des familles et représentants}.
\]

Le tableau de bord n'a pas vocation à remplacer la page Signaux. Il agit comme
un niveau de synthèse qui oriente l'utilisateur vers les cas prioritaires. La
navigation par symbole ouvre ensuite la page \texttt{/signals} avec le titre et
l'horizon concernés, ce qui permet d'examiner les contributions détaillées.

En ce sens, le tableau de bord joue un rôle charnière : il transforme les
sorties du moteur de signaux en information hiérarchisée, tout en conservant un
accès direct au détail nécessaire à la justification analytique.

\section{Présence des niveaux techniques complémentaires}
Au-delà du score et de l'étiquette, le tableau de bord affiche également certains
niveaux techniques, comme les références de support, de clôture et de résistance.
Ces informations ne modifient pas la logique du moteur de signaux, mais elles
renforcent l'utilité opérationnelle de l'écran en donnant des repères de lecture
immédiats.

Ce point est important du point de vue du mémoire : la restitution ne doit pas
être seulement correcte au sens mathématique ; elle doit aussi être utile. Un
score global sans repères contextuels resterait abstrait. L'ajout de niveaux
techniques contribue à relier le signal quantitatif à la lecture quotidienne d'un
graphique ou d'une fiche titre.

\section{Apport décisionnel du tableau de bord}
Le tableau de bord apporte finalement quatre bénéfices.
\begin{enumerate}
\item Il transforme une structure interne complexe en lecture exploitable par
symbole.
\item Il permet une comparaison transversale des titres au sein d'un horizon.
\item Il rend visible la dimension multi-horizons du moteur de signaux.
\item Il conserve un chemin direct vers l'explication analytique détaillée.
\end{enumerate}

En ce sens, le tableau de bord ne constitue pas un simple affichage final. Il est
la couche de synthèse du système de signaux, c'est-à-dire l'interface où la
méthode devient réellement actionnable pour l'utilisateur final.
